\documentclass[9pt]{article}

\usepackage[margin=0.45in,landscape]{geometry}
\usepackage{multicol}
\usepackage{amsmath,amssymb}
\usepackage{array}
\usepackage{enumitem}

\setlist{noitemsep}

\begin{document}

\begin{center}
\Large \textbf{CSCE 222 Exam 1 Cheat Sheet}
\end{center}

\begin{multicols}{2}

%%%%%%%%%%%%%%%%%%%%%%%%%%%%%%%%%%%%%%%%%%%%%%%%
\section*{Logic}

\textbf{Proposition}: statement with truth value.

Variables:

\[
p,q,r
\]

Truth values:

\[
T,F
\]

Logical operators

\begin{tabular}{c|c}
Symbol & Meaning\\
\hline
$\neg p$ & NOT\\
$p\land q$ & AND\\
$p\lor q$ & OR\\
$p\rightarrow q$ & implication\\
$p\leftrightarrow q$ & biconditional
\end{tabular}

Implication equivalence

\[
p\rightarrow q \equiv \neg p \lor q
\]

%%%%%%%%%%%%%%%%%%%%%%%%%%%%%%%%%%%%%%%%%%%%%%%%
\section*{Logical Equivalences}

\textbf{Identity}

\[
p\land T=p
\]

\[
p\lor F=p
\]

\textbf{Domination}

\[
p\lor T=T
\]

\[
p\land F=F
\]

\textbf{Idempotent}

\[
p\lor p=p
\]

\[
p\land p=p
\]

\textbf{Double Negation}

\[
\neg(\neg p)=p
\]

\textbf{De Morgan}

\[
\neg(p\land q)=\neg p\lor\neg q
\]

\[
\neg(p\lor q)=\neg p\land\neg q
\]

%%%%%%%%%%%%%%%%%%%%%%%%%%%%%%%%%%%%%%%%%%%%%%%%
\section*{Rules of Inference}

\begin{tabular}{|c|c|}
\hline
Form & Name\\
\hline
$p, p\rightarrow q \Rightarrow q$ & Modus Ponens\\
\hline
$\neg q, p\rightarrow q \Rightarrow \neg p$ & Modus Tollens\\
\hline
$p\rightarrow q, q\rightarrow r \Rightarrow p\rightarrow r$ & Hypothetical Syllogism\\
\hline
$p\lor q, \neg p \Rightarrow q$ & Disjunctive Syllogism\\
\hline
$p \Rightarrow p\lor q$ & Addition\\
\hline
$p\land q \Rightarrow p$ & Simplification\\
\hline
$p,q \Rightarrow p\land q$ & Conjunction\\
\hline
$p\lor q,\neg p\lor r \Rightarrow q\lor r$ & Resolution\\
\hline
\end{tabular}

%%%%%%%%%%%%%%%%%%%%%%%%%%%%%%%%%%%%%%%%%%%%%%%%
\section*{Logical Fallacies}

\begin{tabular}{|c|c|}
\hline
Form & Error\\
\hline
$p\rightarrow q$

$q$

$\therefore p$
& Affirming Consequent\\
\hline

$p\rightarrow q$

$\neg p$

$\therefore \neg q$
& Denying Antecedent\\
\hline
\end{tabular}

%%%%%%%%%%%%%%%%%%%%%%%%%%%%%%%%%%%%%%%%%%%%%%%%
\section*{Predicate Logic}

Predicate:

\[
P(x)
\]

Universal:

\[
\forall x P(x)
\]

Existential:

\[
\exists x P(x)
\]

Negation rules

\[
\neg(\forall xP(x))\equiv \exists x\neg P(x)
\]

\[
\neg(\exists xP(x))\equiv \forall x\neg P(x)
\]

%%%%%%%%%%%%%%%%%%%%%%%%%%%%%%%%%%%%%%%%%%%%%%%%
\section*{Proof Techniques}

\textbf{Direct}

Assume premise

derive conclusion.

\textbf{Contrapositive}

\[
P\rightarrow Q
\]

prove

\[
\neg Q\rightarrow \neg P
\]

\textbf{Contradiction}

assume opposite

derive impossible statement.

%%%%%%%%%%%%%%%%%%%%%%%%%%%%%%%%%%%%%%%%%%%%%%%%
\section*{Proof by Contradiction Example}

Claim:

\[
\sqrt2\notin \mathbb Q
\]

Assume

\[
\sqrt2=\frac ab
\]

with $gcd(a,b)=1$

Then

\[
2b^2=a^2
\]

$a$ even

$a=2k$

\[
b^2=2k^2
\]

$b$ even

contradiction.

%%%%%%%%%%%%%%%%%%%%%%%%%%%%%%%%%%%%%%%%%%%%%%%%
\section*{Proof by Cases Example}

Prove

\[
n^2+n \text{ even}
\]

Case 1:

\[
n=2k
\]

Case 2:

\[
n=2k+1
\]

Both yield even result.

%%%%%%%%%%%%%%%%%%%%%%%%%%%%%%%%%%%%%%%%%%%%%%%%
\section*{Even/Odd}

Even

\[
n=2k
\]

Odd

\[
n=2k+1
\]

for

\[
k\in \mathbb Z
\]

%%%%%%%%%%%%%%%%%%%%%%%%%%%%%%%%%%%%%%%%%%%%%%%%
\section*{Rational/Irrational}

Rational

\[
x=\frac ab
\]

\[
a,b\in\mathbb Z
\]

Irrational

cannot be expressed as ratio.

%%%%%%%%%%%%%%%%%%%%%%%%%%%%%%%%%%%%%%%%%%%%%%%%
\section*{Sets}

Subset

\[
A\subseteq B
\]

iff

\[
\forall x(x\in A\rightarrow x\in B)
\]

Union

\[
A\cup B
\]

Intersection

\[
A\cap B
\]

Difference

\[
A-B
\]

%%%%%%%%%%%%%%%%%%%%%%%%%%%%%%%%%%%%%%%%%%%%%%%%
\section*{Set Identities}

\begin{tabular}{|c|c|}
\hline
Identity & Name\\
\hline
$A\cup \emptyset=A$ & Identity\\
$A\cap U=A$ & Identity\\
\hline
$A\cup U=U$ & Domination\\
$A\cap\emptyset=\emptyset$ & Domination\\
\hline
$A\cup A=A$ & Idempotent\\
$A\cap A=A$ & Idempotent\\
\hline
$A\cup B=B\cup A$ & Commutative\\
$A\cap B=B\cap A$ & Commutative\\
\hline
$A\cup(B\cup C)=(A\cup B)\cup C$ & Associative\\
\hline
$A\cup(B\cap C)=(A\cup B)\cap(A\cup C)$ & Distributive\\
\hline
$(A\cap B)^c=A^c\cup B^c$ & DeMorgan\\
$(A\cup B)^c=A^c\cap B^c$ & DeMorgan\\
\hline
$A\cup A^c=U$ & Complement\\
$A\cap A^c=\emptyset$ & Complement\\
\hline
\end{tabular}

%%%%%%%%%%%%%%%%%%%%%%%%%%%%%%%%%%%%%%%%%%%%%%%%
\section*{Cardinality}

Cardinality

\[
|A|
\]

Power set

\[
|P(A)|=2^{|A|}
\]

Cartesian product

\[
|A\times B|=|A||B|
\]

%%%%%%%%%%%%%%%%%%%%%%%%%%%%%%%%%%%%%%%%%%%%%%%%
\section*{Functions}

Function

\[
f:A\to B
\]

Domain = $A$

Codomain = $B$

Range

\[
\{f(x)|x\in A\}
\]

%%%%%%%%%%%%%%%%%%%%%%%%%%%%%%%%%%%%%%%%%%%%%%%%
\section*{Function Types (Logical Meaning)}

Injective

\[
\forall x_1,x_2\in A
\]

\[
f(x_1)=f(x_2)\Rightarrow x_1=x_2
\]

Surjective

\[
\forall y\in B,\exists x\in A
\]

\[
f(x)=y
\]

Bijective

injective AND surjective.

%%%%%%%%%%%%%%%%%%%%%%%%%%%%%%%%%%%%%%%%%%%%%%%%
\section*{How to Prove Function Properties}

\textbf{Injective}

1 Assume

\[
f(x_1)=f(x_2)
\]

2 substitute definition

3 show

\[
x_1=x_2
\]

\textbf{Disprove Injective}

Find

\[
x_1\ne x_2
\]

such that

\[
f(x_1)=f(x_2)
\]

%%%%%%%%%%%%%%%%%%%%%%%%%%%%%%%%%%%%%%%%%%%%%%%%

\textbf{Surjective}

1 Let

\[
y\in B
\]

2 solve

\[
y=f(x)
\]

3 show solution $x\in A$

\textbf{Disprove Surjective}

Find

\[
y\in B
\]

such that

\[
\not\exists x
\]

with

\[
f(x)=y
\]

%%%%%%%%%%%%%%%%%%%%%%%%%%%%%%%%%%%%%%%%%%%%%%%%
\section*{Algorithms}

Common complexities

\begin{tabular}{c|c}
Type & Order\\
\hline
Constant & $O(1)$\\
Logarithmic & $O(\log n)$\\
Linear & $O(n)$\\
Quadratic & $O(n^2)$\\
Exponential & $O(2^n)$
\end{tabular}

%%%%%%%%%%%%%%%%%%%%%%%%%%%%%%%%%%%%%%%%%%%%%%%%
\section*{Asymptotic Notation}

Big O

\[
f(n)\le Cg(n)
\]

Big Omega

\[
f(n)\ge Cg(n)
\]

Big Theta

\[
Cg(n)\le f(n)\le C_2g(n)
\]

%%%%%%%%%%%%%%%%%%%%%%%%%%%%%%%%%%%%%%%%%%%%%%%%

%%%%%%%%%%%%%%%%%%%%%%%%%%%%%%%%%%%%%%%%%%%%%%%%
%%%%%%%%%%%%%%%%%%%%%%%%%%%%%%%%%%%%%%%%%%%%%%%%
\section*{Prime and Composite Numbers (Logical Definitions)}

\textbf{Prime Number}

An integer $n$ is \textbf{prime} if

\[
n \text{ is prime} \iff (n>1) \land 
(\forall r,s \in \mathbb{Z}^+)
\Big((n = rs) \rightarrow 
\big((r = 1 \land s = n) \lor (s = 1 \land r = n)\big)\Big)
\]

Meaning:

\begin{itemize}
\item $n>1$
\item whenever $n$ is written as a product $n = rs$ of positive integers,
the only possibilities are
\[
1\cdot n \quad \text{or} \quad n\cdot 1
\]
\end{itemize}

%%%%%%%%%%%%%%%%%%%%%%%%%%%%%%%%%%%%%%%%%%%%%%%%

\textbf{Composite Number}

An integer $n$ is \textbf{composite} if

\[
n \text{ is composite} \iff (n>1) \land
(\exists r,s \in \mathbb{Z}^+)
\Big((n = rs) \land 
\neg\big((r = 1 \land s = n) \lor (s = 1 \land r = n)\big)\Big)
\]

Equivalent useful form:

\[
n>1 \land (\exists r,s \in \mathbb{Z}^+)
(n = rs \land 1<r<n \land 1<s<n)
\]

Meaning:

$n$ can be written as a product of two smaller positive integers.

%%%%%%%%%%%%%%%%%%%%%%%%%%%%%%%%%%%%%%%%%%%%%%%%
%%%%%%%%%%%%%%%%%%%%%%%%%%%%%%%%%%%%%%%%%%%%%%%%
\subsection*{Divisibility}

$a$ divides $b$

\[
a \mid b
\]

means

\[
b = ak
\]

for some

\[
k \in \mathbb{Z}
\]

Example

\[
3 \mid 12
\]

since

\[
12 = 3\cdot4
\]

%%%%%%%%%%%%%%%%%%%%%%%%%%%%%%%%%%%%%%%%%%%%%%%%
\subsection*{Fundamental Theorem of Arithmetic}

Every integer

\[
n>1
\]

can be written as a product of primes:

\[
n = p_1^{e_1}p_2^{e_2}\cdots p_k^{e_k}
\]

and this factorization is \textbf{unique} (up to order).

Example

\[
60 = 2^2 \cdot 3 \cdot 5
\]

%%%%%%%%%%%%%%%%%%%%%%%%%%%%%%%%%%%%%%%%%%%%%%%%
\subsection*{Useful Proof Facts}

If

\[
p \text{ prime and } p\mid ab
\]

then

\[
p\mid a \quad \text{or} \quad p\mid b
\]

This is often used in contradiction proofs.

%%%%%%%%%%%%%%%%%%%%%%%%%%%%%%%%%%%%%%%%%%%%%%%%
\subsection*{Common Proof Pattern}

To prove a number is composite:

Show

\[
n = ab
\]

where

\[
1<a<n
\]

Example

\[
n^2-n = n(n-1)
\]

For

\[
n>2
\]

both factors are greater than 1

so the number is composite.

%%%%%%%%%%%%%%%%%%%%%%%%%%%%%%%%%%%%%%%%%%%%%%%%
%%%%%%%%%%%%%%%%%%%%%%%%%%%%%%%%%%%%%%%%%%%%%%%%
\subsection*{Example Proof: $n^2 - n$ is Composite}

Claim:

For every integer $n>2$, the number $n^2-n$ is composite.

Proof:

Assume $n>2$ and let

\[
m = n^2-n
\]

Factor the expression:

\[
m = n^2-n = n(n-1)
\]

Let

\[
r = n-1, \quad s = n
\]

Then

\[
m = rs
\]

Since $n>2$, we have

\[
1 < n-1 < n^2-n
\]

and

\[
1 < n < n^2-n
\]

Thus

\[
1 < r < m \quad \text{and} \quad 1 < s < m
\]

Therefore

\[
(\exists r,s \in \mathbb{Z}^+)
(m = rs \land 1 < r < m \land 1 < s < m)
\]

By the definition of composite numbers,

\[
n^2 - n \text{ is composite.}
\]

\[
\blacksquare
\]
%%%%%%%%%%%%%%%%%%%%%%%%%%%%%%%%%%%%%%%%%%%%%%%%
\end{multicols}

\end{document}